\subsection{CP-01: operación del recinto desconectado}
\label{caso:cp-01}
\textbf{Trazabilidad:} C07 RT-03.10; SD4/SD5; \texttt{1.10.2.e.5}. Datos preparará el manifiesto inicial de embarcaciones, amarras, permisos, escuela, stock y registros vigentes contra su autoridad, además de operaciones nuevas para cada función local exigible. Se cortarán todos los enlaces exteriores durante 24 horas completas, manteniendo la carga y operación autorizadas del recinto; Calidad comprobará recepción y asignación de embarcaciones, salida/regreso, zarpe, combustible y escuela. Se exigirá ejecución electrónica durante todo el horizonte, persistencia durable, cero efectos duplicados y consulta coherente de autoridad. El diario manual se registrará como contingencia y no como función electrónica cumplida. Evidencia: maniobra de corte, energía/conectividad, registros temporales y conciliación por función.

\subsection{CP-02: terreno y boyas fuera de cobertura}
\label{caso:cp-02}
\textbf{Trazabilidad:} C07 RT-03.10; \texttt{1.10.2.e.5}. Se prepararán dispositivos autorizados de pantalán, varadero, escuela y embarcación de inspección con conjuntos propios del perfil y registro de batería. Se ejecutarán tareas y hechos nuevos durante 12 horas continuas de terreno y ocho horas de inspección, incluyendo reinicio del dispositivo y almacenamiento de evidencia autorizado. Se exigirá conservación de cada registro y su identidad, cero duplicados al reconectar y protección de datos locales. La capacidad se comprobará sobre la jornada completa; un dispositivo con cola vacía no acredita el horizonte. Se registrará separadamente cualquier función que requiera autoridad compartida para no equiparar captura local con decisión global disponible.

\subsection{CP-03: sincronización después de 24 horas}
\label{caso:cp-03}
\textbf{Trazabilidad:} C07 RT-03.13; SD5 5.4; \texttt{1.10.2.e.5}. Se utilizará el universo acumulado en CP-01, con nuevas operaciones concurrentes, duplicados de transporte y mensajes fuera de orden. Tras restituir el enlace, se medirá desde reconexión hasta recepción, aplicación durable, actualización y conciliación completa. Se exigirá duración máxima de treinta minutos, cero pérdida de salidas, regresos, despachos o asistencia y cero efectos repetidos. El manifiesto incluirá tamaño, cantidad, caudal efectivo y hechos nuevos; una estimación de transferencia no sustituirá el tiempo total observado. Datos y Calidad comprobarán cada ID y diferencias; no se purgará la cola para obtener el umbral.

\subsection{CP-04: situación escolar y permisos revocados}
\label{caso:cp-04}
\textbf{Trazabilidad:} C07 RT-05.29; SD5, autoridad escolar. Se prepararán alumnos, apoderados, monitores, botes y permisos vigentes, vencidos y revocados, con datos protegidos. Se registrarán salida, cambio, regreso y retiro desde dos puestos autorizados, incluyendo partición de comunicación y reintento. Se exigirá reflejo de cada hecho contra la autoridad compartida en tiempo real, rechazo de actuación sin autorización y cero acceso cruzado entre apoderados. Una partición que impida conocer quién está en el agua se registrará como incumplimiento aunque exista copia local. Evidencia: secuencia de hechos, consulta de ambos puestos y decisiones de permisos.

\subsection{CP-05: tiempos completos de operación}
\label{caso:cp-05}
\textbf{Trazabilidad:} C07 RT-09.01; \texttt{1.10.2.e.4/10}. Funcional preparará embarcación visitante compatible y no compatible, registros de salida/regreso, personas y antecedentes de zarpe y clase de hasta 18 botes. Usuarios representativos ejecutarán las tareas con dispositivos y condiciones de uso seguras del perfil. Se medirán consulta radial a confirmación de amarra $\leq45$ s; inicio a persistencia de salida/regreso $\leq10$ s y máximo dos interacciones; nómina/legajo $\leq90$ s; y clase completa $\leq60$ s. Se conservarán tiempos individuales y errores; medir sólo API no acredita la tarea. La aceptación exigirá resultados correctos dentro de cada máximo.

\subsection{CP-06: avisos de emergencia y planes vencidos}
\label{caso:cp-06}
\textbf{Trazabilidad:} C07 RT-09.01 y RT-16.21; módulo ALE. Se prepararán 296 destinatarios de ensayo con mensaje de texto, mensajería instantánea y correo, contactos válidos, errores de entrega y acuses ausentes, más planes sin cierre con vencimiento identificado. Se activará una emergencia y se comprobarán los tres canales simultáneos, reintentos y escalamiento a llamada registrada para quienes no acusen; los destinos de ensayo evitarán avisos reales no autorizados. El último envío deberá ocurrir dentro de diez minutos, registrando aparte entregas y acuses por destinatario. La alerta de plan vencido se generará en máximo quince minutos desde vencimiento, sin agregar otro período de gracia. Evidencia: destinatarios, intentos, respuestas, pendientes y escalamiento.

\subsection{CP-07: combustible y efectos externos idempotentes}
\label{caso:cp-07}
\textbf{Trazabilidad:} SD5, autoridad local y transacciones; C07 RT-03.10/13. Se prepararán operaciones de despacho y hechos facturables con ID y versión, confirmaciones repetidas y respuestas externas inciertas. Se repetirán solicitudes, se interrumpirá comunicación después de persistir y antes del acuse y se reiniciará el productor. Se exigirá un único efecto por hecho y correspondencia entre movimiento, evidencia y resultado externo. Una respuesta incierta se consultará y conciliará según el contrato; no se reintentará como fallo definitivo si eso pudiera duplicar un cargo. Evidencia: IDs, transacciones, inbox/outbox, acuses y conciliación.

\subsection{CP-08: migración y retorno con hechos nuevos}
\label{caso:cp-08}
\textbf{Trazabilidad:} BT RT-05.13/14, cap.~20; \texttt{1.10.1.e.j.k}. Datos preparará el universo histórico exigible de SD5, incluidos los catorce expedientes de abandono, relaciones y documentos asociados, con manifiesto protegido y reglas versionadas. Cada ensayo recorrerá extracción, transformación, carga, conciliación y retorno, incluyendo hechos nuevos posteriores al corte y resultados externos inciertos. Se exigirá cero diferencias no explicadas y las tolerancias críticas de SD5, sin compensar saldos de signo contrario ni descartar hechos posteriores al restaurar. Ambos ensayos deberán ser completos y válidos para la versión definitiva; el segundo incluirá suplentes técnico y funcional. Evidencia: manifiestos, tiempos por fase, diferencias y decisión de retorno.

\subsection{CP-09: conmutación regional y retorno DR-4}
\label{caso:cp-09}
\textbf{Trazabilidad:} BT RT-07.01--06; \texttt{1.10.2.e.7} y \texttt{1.12.4}. Operación preparará versión, réplica, identidades, claves, objetos, colas y corte de hechos confirmados en los sitios comprometidos. Se aislará el origen, se promoverá el destino y se comprobará capacidad, autenticación, datos y operaciones antes del retorno. Se exigirá RTO $\leq4$ h desde interrupción hasta funcionalidad restituida y RPO $\leq15$ min sobre hechos regionales confirmados, sin doble escritor ni efectos repetidos. Se distinguirán hechos de terreno aún no recibidos y su obligación local. Evidencia: cronología, último dato confirmado, permisos, fencing, conciliación y retorno; una réplica existente no acredita el ensayo.

\subsection{CP-10: acceso cruzado y revocación}
\label{caso:cp-10}
\textbf{Trazabilidad:} BT cap.~11; SD4, identidad y segregación; \texttt{1.10.2.e.8}. Seguridad preparará cuentas de los perfiles comprometidos, sesiones emitidas, credenciales revocadas y solicitudes a recursos propios y ajenos. Se ejecutarán accesos negativos por interfaz directa y portal, revocación durante sesión y comprobación de registros sin exposición innecesaria de información personal. Se exigirá cero lectura/escritura no autorizada y cierre de todos los hallazgos críticos/altos. La intrusión independiente conservará alcance y reporte íntegro; el análisis automático complementa esa comprobación y no la reemplaza.

\subsection{CP-11: accesibilidad y aprendizaje por perfil}
\label{caso:cp-11}
\textbf{Trazabilidad:} BT RT-13.01--11; \texttt{1.10.2.e.9/10}. Se prepararán tareas por perfil, dispositivos soportados, navegación de teclado y ayudas técnicas, además de sesiones de instrucción registradas. Se comprobarán criterios WCAG 2.2 AA aplicables automática y manualmente, foco, errores, objetivos táctiles y uso seguro con guantes. Se medirán tiempo, interacciones, ayuda, errores y aprendizaje con las fórmulas de SD9, sin promediar perfiles ni borrar recuperaciones. Se exigirá 95\,\% de tareas sin asistencia, error $\leq5\,\%$, actuaciones críticas 100\,\% correctas y objetivo de interfaz tras cuatro horas de instrucción; la formación adicional se registrará como recuperación.

\subsection{CP-12: promoción y reversa sin interrupción}
\label{caso:cp-12}
\textbf{Trazabilidad:} BT RT-04.06/07/10; DEP-6/MIG-6; SD6. Se prepararán artefactos anterior/nuevo, digest, configuración, esquema compatible y solicitudes/conexiones en vuelo sobre EQ-4 a $1,5$ veces peak. Se ejecutarán promoción y retorno automatizados, con falla inducida y hechos nuevos durante el cambio. Se exigirá cero interrupción, pérdida o duplicación, cumplimiento de tiempos y conservación de autoridad y sesión. El ensayo comprobará expandir/transitar/retirar esquema y retorno representable. Una conmutación de infraestructura cuya continuidad no esté demostrada permanecerá bloqueada; una ventana o RTO conforme no dispensará ese criterio.

\subsection{CP-13: crecimiento por parametrización}
\label{caso:cp-13}
\textbf{Trazabilidad:} C07 RT-02.12; SD4 y modelo de SD5. Se preparará una copia de ensayo del catálogo con pantalanes, amarras, puestos de invernada y boyas, manteniendo sus relaciones y reglas de compatibilidad. Una persona usuaria autorizada del perfil de administración configurará nuevos elementos hasta el escenario de 380 amarras y 70 boyas, usando funciones ordinarias y documentación del producto. Se exigirá cero cambio de código o rediseño, cero intervención de especialista para parametrizar y funcionamiento correcto de asignación, búsqueda, permisos, telemetría y consultas sobre las unidades nuevas. El incremento de recursos de capacidad se distinguirá del acto de parametrización y deberá conservar el plan de SD4. Evidencia: parámetros y auditoría del cambio, relaciones, resultados y recursos utilizados.

\subsection{CP-14: cobertura, segregación y enlace de respaldo}
\label{caso:cp-14}
\textbf{Trazabilidad:} C07 RT-03.24. Se preparará el protocolo de recorrido y puntos de medición de los cinco pantalanes y el varadero, con el rango de marea real, flexión de cableado y condiciones seguras de desplazamiento. Se comprobará captura y consulta en cada zona, segregación efectiva entre socios, administración, cámaras y operación, y conmutación al enlace de respaldo contratado. Se exigirá cobertura funcional de todas las zonas, rechazo de accesos no autorizados entre segmentos y continuidad conforme al requisito del servicio durante la falla del enlace principal. El expediente incluirá mediciones, diagrama real, pruebas negativas y evidencia del contrato y recepción del respaldo, sin atribuirlo al equipamiento existente. Una visita en una única condición de marea no cerrará la comprobación del rango.

\subsection{CP-15: retención, custodia y eliminación}
\label{caso:cp-15}
\textbf{Trazabilidad:} C07 RT-05.10; SD5. Se prepararán registros de cada categoría con fechas de origen y vencimiento, originales y copias, bloqueos legales de eliminación y permisos de acceso. El ensayo verificará con tiempo controlado los límites: menores hasta 18 años más cinco y mínimo diez años; salida/regreso y zarpe cinco años; contratos/custodia/cuenta seis; residuos/derrame seis; consumos cinco; fondeo vida útil más cinco; combustible cinco; video seis meses. Para menores se aplicará el vencimiento más tardío entre ambas condiciones. Antes de vencer se exigirá recuperación íntegra por perfil autorizado; después se comprobará el proceso de eliminación segura de las copias bajo su regla y el registro de la actuación, sin borrar evidencia sujeta a custodia legal. Una simulación de fecha verificará la regla, sin declararse una retención histórica ya ejecutada.

\subsection{CP-16: intercambio sectorial y disponibilidad efectiva}
\label{caso:cp-16}
\textbf{Trazabilidad:} C07 RT-05.23. Integración preparará por autoridad o servicio el formato, versión, canal, credenciales, disponibilidad confirmada y contrato de intercambio para zarpe, avisos meteorológicos/marejada, residuos peligrosos, indicadores ambientales y mensajería de baja capacidad. Se comprobarán mensajes válidos, rechazo de formato incorrecto, reintentos, latencia de horas y entrega desordenada cuando correspondan. Se exigirá correspondencia con el formato y canal realmente admitidos y persistencia/conciliación del resultado. Un simulador validará el desarrollo, pero no cerrará la disponibilidad efectiva del tercero. Si una autoridad no dispone de API, se comprobará el procedimiento seleccionado y autorizado de intercambio, sin inventar interfaz o aceptación institucional.

\subsection{CP-17: actualización de datos y mezcla de peak}
\label{caso:cp-17}
\textbf{Trazabilidad:} C07 RT-05.29 y RT-09.02; SD5 5.4. Se preparará la población y mezcla derivada de la volumetría 14.1/14.2, con mañana de fin de semana largo, consultas, avisos, telemetría, historia retenida y hechos nuevos simultáneos. El ensayo a $1,5$ veces peak conservará los denominadores y tasas de SD4--SD5. Se verificará ocupación en máximo 60 s desde evento, acumulado por amarra al menos horario, saldo y situación escolar en tiempo real contra su autoridad y consolidación ambiental diaria. El informe conservará P95 de respuesta y máximos específicos del caso, cargas, recursos, errores y saturación. Una consulta sobre proyección antigua deberá indicar su condición sin presentarse como dato actual; un promedio anual no sustituirá la concentración del peak.

\subsection{CP-18: autonomía administrada y periféricos}
\label{caso:cp-18}
\textbf{Trazabilidad:} C07 RT-06.01 y RT-17.06; SD4/T-11. Se preparará el inventario seleccionado, variantes, emplazamiento y certificados aplicables de gabinetes, medidores, sensores, lectores, despacho de combustible, estación meteorológica, control de acceso y las 38 cámaras existentes. Se comprobarán compatibilidad e integración efectiva, unidades y calibración cuando corresponda, recepción de avisos y comportamiento ante fallas. Los gabinetes nuevos de administración, pantalán y varadero acreditarán la protección exigible; la instrumentación del surtidor conservará la aptitud para su clasificación. Durante operación normal se comprobará autoadministración y ausencia de intervención humana local, distinguiendo mantenimiento o incidentes extraordinarios. Un catálogo no acredita instalación ni recepción y una prueba de escritorio no reemplazará la verificación en torreta flotante.

\subsection{CP-19: personas, perfiles y bilingüismo}
\label{caso:cp-19}
\textbf{Trazabilidad:} C07 RT-12.11/12, RT-13.08/12, RT-17.01 y RT-22.04. Se preparará la matriz de perfiles internos/externos, incluyendo armadores, socios, apoderados, visitantes, contratistas/trabajadores, charter, proveedores/gestores y autoridades aplicables. Los cinco perfiles móviles recorrerán sus tareas: pantalán, escuela, varadero, inspección y armador en baja capacidad. Se comprobarán guantes, pantalla mojada y sol directo en pruebas seguras; no se exigirá interacción durante maniobra con carga suspendida ni con manos ocupadas en una actuación insegura. Visitantes, tripulaciones extranjeras y charter usarán todas sus interfaces y comunicaciones en español e inglés. La formación y evaluación incluirán personal nuevo y monitores renovados antes del congelamiento, impartidas por Synaptix, sin depender de un administrador o formador del CLIENTE. Evidencia: tareas/perfiles, idioma, condiciones, incidentes de uso y evaluación individual.

\subsection{CP-20: privacidad, auditoría y firma}
\label{caso:cp-20}
\textbf{Trazabilidad:} C07 RT-11.10, RT-16.09/14. Se prepararán datos protegidos de salud/menores, armadores, tripulantes, contratistas y posición consentida, con autorizaciones y revocaciones. Se exigirá auditoría de cada acceso a salud/menores, posición compartida y deuda, además de modificaciones, conservando quién, qué, cuándo y resultado sin exponer innecesariamente el contenido sensible. Los casos de firma cubrirán contrato/renovación de amarra, autorización del apoderado, recepción/entrega de custodia, acta de izaje y entrega de residuos al gestor. Cada modalidad deberá corresponder a la normativa aplicable y comprobar identidad, integridad, vigencia y conservación. No se aceptará un clic genérico como demostración de firma jurídicamente admisible en todos los casos.

\subsection{CP-21: portales, baja capacidad y notificaciones}
\label{caso:cp-21}
\textbf{Trazabilidad:} C07 RT-16.21/30 y RT-17.01. Se prepararán datos del portal público, armador, apoderado, visitante y contratista y destinos de prueba. Sin autenticación se comprobarán disponibilidad/tarifas del servicio de visita, condiciones de acceso, estado y avisos, en español e inglés. Autenticado se verificarán reserva/pago, contrato/consumo/cuenta/documentos/plan del armador, autorización y salida/regreso escolar, acreditación y residuos del contratista. Se comprobarán avisos de retorno a apoderados y de eventos/cuenta a armadores, más canal de baja capacidad con latencia de horas y mensajes desordenados. Se exigirá correspondencia con la autoridad, aislamiento entre cuentas y un único efecto por solicitud, conservando recibos y estados de entrega.

\subsection{CP-22: servicio gestionado y restricciones de intervención}
\label{caso:cp-22}
\textbf{Trazabilidad:} C07 RT-10.05, RT-15.02, RT-21.06/16 y RT-22.04. Proyecto preparará el calendario real de campañas y 14 eventos, congelamiento del 15 de diciembre al 15 de marzo, avisos de marejada y acceso al campo de 40 boyas; verificará logística marítima segura considerando 6 millas náuticas y 22 km por camino, sin suponer acceso permanente o personal residente. Se comprobarán suspensión inmediata, reprogramación y capacidad de contingencia sin desplazar hitos. El expediente de capacidad incluirá experiencia acreditada en normativa marítima/residuos y servicio gestionado sin TI interna. Se verificarán turnos efectivos de mesa 06:00--24:00 en temporada y 08:00--20:00 el resto, atención en ambos idiomas y emergencia/planes vencidos 24x7x365, con escalamiento y respuesta crítica preservados. Los contratos/cuadrantes y comprobaciones del servicio acreditarán cobertura; un calendario declarado no demuestra personal contratado ni respuesta obtenida.
